\clearpage
\section{Thermal noise: fluctuations and dissipation}
\subsection{Thermal equilibrium and equipartition}
Thermal equilibrium means that macroscopic averages no longer drift, not that microscopic charges stop moving. A voltage can have zero ensemble mean and nonzero mean square even after the circuit has fully settled. The equilibrium voltage distribution is determined by the stored energy and temperature.\footnote{Here $T$ is absolute temperature; observation duration is denoted $T_{\mathrm{obs}}$ when needed.}

\paragraph{From stored energy to voltage probability.}
A capacitor in thermal contact with its surroundings continually gains and loses small amounts of energy. At equilibrium, the statistical weight of a state with energy $U$ is proportional to
\[
e^{-U/(kT)}.
\]
This is the \emph{Boltzmann factor}. It compares the energy cost $U$ with the thermal energy scale $kT$, where $k=1.380649\times10^{-23}\,\mathrm{J/K}$ is Boltzmann's constant. We use this statistical-mechanical result without deriving it here.

For an unbiased linear capacitor, $U_C=\tfrac12 Cv^2$. Substitution gives
\[
p(v)\propto e^{-Cv^2/(2kT)}.
\]
Compare this with the zero-mean Gaussian density,
\[
p(v)=\frac1{\sqrt{2\pi\sigma_v^2}}e^{-v^2/(2\sigma_v^2)}.
\]
Matching the Gaussian exponent gives
\begin{equation}
\boxed{\sigma_v^2=\E[v^2]=\frac{kT}{C},\qquad
\E[U_C]=\frac12 C\E[v^2]=\frac12 kT.}
\label{eq:ktc-equip}
\end{equation}
The quadratic energy gives a Gaussian voltage distribution: a larger $C$ makes a given fluctuation more costly, whereas higher temperature makes it more likely. Equation~\eqref{eq:ktc-equip} is an instance of classical \emph{equipartition}: each independent quadratic energy term has mean energy $kT/2$.

A resistor permits charge and energy exchange. It controls the rate of relaxation, but not the final $kT/C$ variance. For a biased linear capacitor, $v$ denotes the fluctuation about its mean. These classical results assume effectively continuous charge; a Gaussian voltage distribution alone does not determine temporal correlation.

\subsection{Johnson--Nyquist noise from equipartition}
Connect a resistor $R$ to a capacitor $C$ at temperature $T$. Represent the resistor's noise by a series voltage $e_R$ with an unknown, constant one-sided PSD $S_0$. The voltage transfer function is
\[
H(f)=\frac1{1+j2\pi fRC},\qquad
S_{v_C}^{(1)}(f)=S_0|H(f)|^2=\frac{S_0}{1+(2\pi fRC)^2}.
\]
\RCFigure

In the ideal white-noise model, its equilibrium variance is
\begin{align*}
\E[v_C^2]
&=\int_0^\infty S_0|H(f)|^2\dd f
=\frac{S_0}{2\pi RC}\int_0^\infty\frac{\dd u}{1+u^2}\\
&=\frac{S_0}{2\pi RC}\frac\pi2
=\frac{S_0}{4RC},\qquad u=2\pi fRC.
\end{align*}
Equipartition requires this to equal $kT/C$. Therefore
\[
\frac{S_0}{4RC}=\frac{kT}{C}
\quad\Longrightarrow\quad
\boxed{S_0=4kTR.}
\]
Any other white-noise level would give the wrong equilibrium capacitor energy. This argument fixes the level \emph{assuming whiteness}; it does not prove that the noise is white or establish the general impedance formula. For a real resistor, the RC corner must lie well below the frequencies where its resistance or thermal-noise density changes appreciably. Extending the flat spectrum to infinity is then an approximation for evaluating the low-frequency result.

The open-circuit density is therefore $4kTR$ V$^2$/Hz. In a flat band of width $B$,
\begin{equation}
v_{n,\rms}=\sqrt{\int_{\mathrm{band}}S_{v,\mathrm{oc}}^{(1)}(f)\,\dd f}
=\sqrt{4kTRB}.\label{eq:resistor-rms}
\end{equation}
The equivalent parallel short-circuit current PSD is $S_{i,\mathrm{sc}}^{(1)}=4kT/R$ A$^2$/Hz. These are two descriptions of the same noisy resistor, not two independent sources to add simultaneously.

The general fluctuation--dissipation relation extends this result. In the classical regime $hf\ll kT$, a passive linear impedance in equilibrium has one-sided open-circuit PSD
\begin{equation}
\boxed{S_{v,\mathrm{oc}}^{(1)}(f)=4kT\operatorname{Re}Z(f),\qquad f>0.}
\label{eq:fdt}
\end{equation}
Here $h$ is Planck's constant. Only the dissipative part contributes. Its general derivation is beyond our scope; the RC argument establishes the resistor level under the stated white-noise approximation.

\subsection{RC noise: correlation and settling}
\label{sec:rc-chain}
The RC calculation above fixes the total equilibrium variance. The circuit dynamics tell us how that variance is distributed in time and frequency. With $S_{e_R}^{(1)}=4kTR$, the capacitor spectrum is Lorentzian: its corner is $f_p=1/(2\pi RC)$, and its area is the $kT/C$ already calculated.

\paragraph{Equilibrium correlation.}
The circuit equation and impulse response are
\[
RC\frac{\dd v_C}{\dd t}+v_C=e_R(t),\qquad
h(t)=\frac{e^{-t/(RC)}}{RC}\,u(t).
\]
Here $u(t)$ is the unit step. The two-sided driving PSD is $2kTR$, so $R_{e_R}(\tau)=2kTR\delta(\tau)$. In equilibrium, the filter acts on the entire past:
\[
v_C(t)=\int_0^\infty h(s)e_R(t-s)\dd s.
\]
Substitute this convolution at both times into the ensemble correlation:
\begin{align*}
R_{v_C}(\tau)
&=\E[v_C(t)v_C(t+\tau)]\\
&=\int_0^\infty\!\int_0^\infty h(s)h(r)
\E[e_R(t-s)e_R(t+\tau-r)]\dd r\dd s\\
&=2kTR\int_0^\infty\!\int_0^\infty h(s)h(r)
\delta(\tau+s-r)\dd r\dd s.
\end{align*}
For $\tau\geq0$, the delta function selects $r=s+\tau\geq0$, giving
\begin{align*}
R_{v_C}(\tau)
&=2kTR\int_0^\infty h(s)h(s+\tau)\dd s\\
&=\frac{2kTR}{(RC)^2}e^{-\tau/(RC)}\frac{RC}{2}
=\frac{kT}{C}e^{-\tau/(RC)}.
\end{align*}
Even symmetry gives
\begin{equation}
\boxed{R_{v_C}(\tau)=\frac{kT}{C}e^{-|\tau|/(RC)}.}
\end{equation}
Thus $RC$ is both the relaxation time and the correlation time. Increasing $R$ slows the fluctuations and narrows the spectrum without changing the equilibrium variance.

If the drive is jointly Gaussian and the circuit is initialized in equilibrium, the output is a stationary Gaussian process. Its exponential covariance satisfies the mean- and correlation-ergodicity conditions in Sections~\ref{sec:mean-proof} and~\ref{sec:corr-ergodicity}. Whiteness alone does not establish Gaussianity; a raw periodogram still needs averaging to reduce its scatter (Section~\ref{sec:psd-measurement}).

\paragraph{Settling and sampled noise.}
If $v_C(0)=0$ in every outcome, the response instead begins at zero:
\[
v_C(t)=\int_0^t h(t-u)e_R(u)\dd u.
\]
Using the delta correlation of the drive gives
\begin{align*}
\Var[v_C(t)]
&=2kTR\int_0^t h(t-u)^2\dd u\\
&=\frac{2kTR}{(RC)^2}\int_0^t e^{-2(t-u)/(RC)}\dd u\\
&=\boxed{\frac{kT}{C}\left(1-e^{-2t/(RC)}\right).}
\end{align*}
This transient is not stationary. The variance approaches $kT/C$ only after settling. If an ideal switch then isolates the capacitor, its sampled charge is retained; the capacitor does not generate a new independent noise sample at every instant. Leakage permits further relaxation.

An ideal capacitor has no dissipative noise source of its own. Its equilibrium uncertainty comes from energy exchange with its surroundings. Taking $R\to\infty$ before allowing the circuit to settle isolates its initial charge; allowing equilibrium first and then disconnecting retains a random equilibrium sample. The order describes different experiments.

\begin{example}[title={Changing resistance at fixed capacitance}]
At $T=290\,\mathrm K$, $C=1\,\mathrm{pF}$, and $R=1\,\mathrm{k\Omega}$,
\[
\sqrt{4kTR}=4.00\,\mathrm{nV}/\sqrt{\mathrm{Hz}},\qquad
f_p=159.2\,\mathrm{MHz},\qquad
v_{C,\rms}=63.3\,\mu\mathrm V.
\]
Replacing $R$ by $10\,\mathrm{k\Omega}$ increases the resistor ASD by $\sqrt{10}$, reduces $f_p$ to $15.92\,\mathrm{MHz}$, and increases the settling time tenfold. The equilibrium RMS voltage remains $63.3\,\mu\mathrm V$: the higher source density is offset by the narrower bandwidth.
\end{example}

\subsection{Available noise power}
Represent a resistor by an ideal noiseless resistor $R$ in series with its open-circuit noise voltage $v_n$. Attach an otherwise noiseless matched load $R_L=R$, as in Figure~\ref{fig:match}. The voltage divider gives $v_L=v_n/2$, so
\begin{align}
S_{v_L}^{(1)}&=\left|\frac{R}{R+R}\right|^2(4kTR)=kTR,\\
P_{n,L}&=\frac{\E[v_L^2]}{R}
=\frac{kTRB}{R}=\boxed{kTB}.\label{eq:ktb}
\end{align}
The divider ratio is squared because PSD tracks mean square. Thus $4kTR$ is an open-circuit voltage PSD, whereas $kT$ is the available power PSD delivered to a matched load. This calculation includes only source noise. Equal-temperature resistors also exchange noise in the reverse direction, so their net equilibrium heat flow is zero.
\MatchedResistorFigure

\Needspace{5\baselineskip}
\begin{example}[title={Thermal noise density at 290 K}]
At $T_0=290\,\mathrm K$, $kT_0=4.0039\times10^{-21}\,\mathrm{W/Hz}$. Decibel-milliwatts are defined by $P_{\mathrm{dBm}}=10\log_{10}(P/1\,\mathrm{mW})$, with $P$ in watts. In a $1\,\mathrm{Hz}$ band,
\[
10\log_{10}\!\left(\frac{kT_0(1\,\mathrm{Hz})}{1\,\mathrm{mW}}\right)
=-173.98\,\mathrm{dBm}.
\]
Thus the conventional available noise-power density is $-173.98\,\mathrm{dBm/Hz}$. Increasing a flat integration band from $1\,\mathrm{MHz}$ to $10\,\mathrm{MHz}$ raises integrated power by $10$ dB and RMS voltage, at fixed resistance, by $\sqrt{10}$.
\end{example}

\subsection{Beyond the classical limit}
A white spectrum cannot extend to arbitrarily high frequency. For one matched traveling-wave mode, the emitted thermal power density is
\[
S_{P,\mathrm{emit}}(f)=\frac{hf}{e^{hf/(kT)}-1},\qquad f>0.
\]
It approaches $kT$ when $hf\ll kT$ and rolls off when $f$ approaches $kT/h$ (about $6.04\,\mathrm{THz}$ at $290\,\mathrm K$). Substituting $x=hf/(kT)$ yields
\[
P_{\mathrm{emit}}=\frac{(kT)^2}{h}\int_0^\infty\frac{x}{e^x-1}\dd x
=\frac{\pi^2(kT)^2}{6h}\simeq39.8\,\mathrm{nW}.
\]
The integral uses the standard identity $\int_0^\infty x/(e^x-1)\dd x=\pi^2/6$. One route is to expand $(e^x-1)^{-1}=\sum_{n\geq1}e^{-nx}$ and integrate term by term: $\int_0^\infty xe^{-nx}\dd x=1/n^2$, whose sum is $\pi^2/6$.

This emitted power applies to a single matched traveling-wave mode. Equal-temperature terminations exchange equal powers. The symmetrized quantum noise spectrum includes a zero-point term and differs from emitted power; a detector's ordering convention must be specified before extending a classical formula into the quantum regime. Parasitic inductance, capacitance, frequency-dependent loss, and the measurement network also constrain real circuit bandwidth.
